\section{Introduction}\label{sec:introduction}

Let $\mathcal P$ denote the set of positive primes, and put
$\N_0=\{0,1,2,\ldots\}$.  For $A,B\subseteq\N_0$, write
$A+B=\{a+b:a\in A,\ b\in B\}$.  Two subsets of $\N_0$ are
\emph{asymptotically equal} if their symmetric difference is finite.
Ostmann's conjecture asserts that $\mathcal P$ is asymptotically
additively indecomposable: it is not asymptotically equal to $A+B$ when
both summands contain at least two elements.  The conjecture goes back
to his 1956 treatise~\cite[p.~13]{Ostmann1956}; the eventual-equality
formulation and its terminology are recorded explicitly in
\cite[Definition~1.1 and Conjecture~1.2]{ElsholtzHarper}.
It is also known as the inverse Goldbach problem.

\begin{theorem}[Ostmann's conjecture]\label{cor:ostmann}
If $A,B\subseteq\N_0$ satisfy $|A|,|B|\ge2$, then
$(A+B)\mathbin{\triangle}\mathcal P$ is infinite.
Equivalently, every finite modification of $\mathcal P$ is additively
indecomposable into two sets with at least two elements each.
\end{theorem}

Laffer and Mann showed that a hypothetical decomposition must have two
infinite summands~\cite[Theorem~12]{LafferMann1964}.  We first give a
short sieve proof of this reduction in Lemma~\ref{lem:finite-factor},
then state the two-infinite-set contradiction as
Theorem~\ref{thm:main}.  The rest of the paper proves that theorem.

The conclusion concerns both eventual coverage of the primes and
eventual exclusion of composite sums.  Neither requirement is replaced
by a density condition.  In particular, if $|A|,|B|\ge2$ and $A+B$
contains every sufficiently large prime, then it contains infinitely
many composite numbers.

Earlier work established increasingly strong restrictions on a
hypothetical decomposition.  Hornfeck's work
\cite{Hornfeck1954,Hornfeck1955} was followed by sieve arguments of
Pomerance, S\'ark\"ozy and Stewart~\cite{PomeranceSarkozyStewart1988},
Hofmann and Wolke~\cite{HofmannWolke1996}, and
Elsholtz~\cite{Elsholtz2001Remark}.  Write $A(x)=|A\cap[0,x]|$ and
similarly for $B$.  Bounds such as $A(x)B(x)\ll x$ are compatible with
the coverage lower bound $A(x)B(x)\gg x/\log x$ and therefore do not
by themselves exclude a decomposition.  Elsholtz combined the large
and larger sieves to put both counting functions at the square-root
scale, up to powers of $\log x$, and ruled out a decomposition into
three nontrivial summands~\cite{Elsholtz2001}.  His later refinement
\cite[Theorem~1.9]{Elsholtz2006} gives the bounds reproduced in
Lemma~\ref{lem:size}.  Croot and Elsholtz also studied thin ternary
sumsets contained in the primes, obtaining restrictions under a
regularity hypothesis on representation multiplicities
\cite[Theorem~1]{CrootElsholtz2005}.  Shao later proved a finite ternary
obstruction: for some absolute $c>0$, three subsets of $[1,N]$, each
of size at least $N^{1/3-c}$, have a composite number in their sumset
once $N$ is sufficiently large
\cite[Theorem~1.3 of the preprint version]{Shao2016}.
In their positive-integer formulation,
Elsholtz and Harper sharpened the binary counting bounds to
\[
 \frac{\sqrt x}{\log x\log\log x}
 \ll A(x),B(x)\ll\sqrt x\log\log x
\]
under a hypothetical eventual decomposition
\cite[Theorem~2.6]{ElsholtzHarper}.

Green and Harper developed inverse questions for the large sieve and
proved that a suitable inverse-sieve conjecture would imply Ostmann's
conjecture~\cite[Conjecture~1.5 and Theorem~1.8]{GreenHarper}.
Their conjecture proposes a quadratic description of sets near the
square-root sieve bound.  Under half-residue restrictions
$|C\bmod p|\le p/2+O(1)$ at every prime and $|C|\gg\sqrt N$ for
$C\subseteq[1,N]$, Hanson proved additive correlation with the squares
and a logarithmic-size intersection with a quadratic image
\cite[Theorem~1.2 and Corollary~1.3]{Hanson2020}.
Croot, Mao and Yip subsequently proved inverse theorems for local
restrictions to short arithmetic progressions, with extensions to unions
of progressions and other additive structures~\cite{CrootMaoYip2025}.
More recently, Croot, Mao, Pohoata and Yip used a weighted entropy
argument to derive a further necessary condition for a hypothetical
decomposition of the primes into two sets of positive integers, each
containing at least two elements
\cite[Theorem~1.10 and Corollary~1.11]{CrootMaoPohoataYip2026}.
For some $c>0$ and every sufficiently large $N$, each summand then has at
least $\exp(c\sqrt{\log N}/(\log\log N)^{3/2})$ elements in $[1,N]$
lying in an integral quadratic image $m\pm\Z^2$; the product of the two
intersection sizes is at least
$\exp(c\sqrt{\log N}/\log\log N)$.  These quadratics may depend on $N$.

These large intersections with quadratic images fall short of the
near-containment required by Green and Harper's conditional route.
Our proof uses the simultaneous residue
restrictions and coverage of every sufficiently large prime to reach a
contradiction without such a classification.

Under a hypothetical decomposition,
put $D=-B$.  For each prime $p$, deleting finite initial segments from
$A$ and $B$ leaves disjoint images of $A$ and $D$ in $\F_p$.
Completing these subsets
to a partition gives a residue set $S_p$ and its complement.  The
collision estimate in Section~\ref{sec:preliminaries} shows, on suitable
prime averages, that the two parts have approximately equal size and
that the summands are approximately uniform on their respective parts.

The next stage rules out persistent correlations with translated
multiplicative characters.  Section~\ref{sec:quadratic} treats quadratic
characters with independently chosen translating centers at every
prime.  A moment argument and Poisson summation force a common rational
center. The quadratic large sieve then confines positive fractions of
long tails of the summands to a small family of quadratic
kernels, leading to a collision contradiction.  Section~\ref{sec:characters} treats all higher character
orders by repeated Cauchy--Schwarz transfers.  Its anchor variables
distinguish the copied terms and supply the character cancellation
needed to control the diagonals.  Together these arguments give the
mixed-character decorrelation in Corollary~\ref{cor:mixed-flatness}.

Prime coverage enters separately in Section~\ref{sec:supply}.  A tensor
estimate shows that the normalized additive transforms of the residue
indicators cannot have small $L^1$ norm on too much harmonic prime
mass.  The proof compares a nonnegative sumset statistic with its
average over the primes.  This comparison retains the coverage
information beyond the preliminary lower bounds for the summands.

The remaining residue indicators need not have a prescribed algebraic
form.  Section~\ref{sec:tree} proves a finite-field comparison for the
binary trees produced by the transfers, using the mixed-character
decorrelation. Its hypotheses involve only a probability $L^2$ bound
and mixed-character correlations; they permit highly nonuniform
pointwise values. This feature is needed for the exact Fourier
transforms of the residue indicators.  Sections~\ref{sec:construction} and~\ref{sec:arithmetic}
construct a positive statistic and establish its comparison estimates.
The arithmetic separation includes repeated internal prime labels,
coprimality conditions, and the possible exceptional real character
in prime progression estimates.  Finally, Section~\ref{sec:conclusion}
averages over permutations of the bulk variables.  Most pairs of arrangements have few overlap components and hence
small correlation. Their total contribution and the remaining diagonal
terms contradict the lower bound inherited from the positive statistic.
The transfer depth is chosen sufficiently large and fixed before the
large scale tends to infinity.

This proves asymptotic indecomposability directly from the simultaneous
residue restrictions and prime coverage.  The argument does not require
the general inverse-sieve conjecture of~\cite{GreenHarper}, and no such
general classification theorem is asserted here.
